\documentclass[12pt,a4paper]{article}
\usepackage[spanish,es-tabla]{babel}
\decimalpoint
\usepackage{array}
\newcolumntype{C}[1]{>{\centering\arraybackslash}m{#1}}
\usepackage{bbm}
\usepackage{bm}
\usepackage{tikz}
\usetikzlibrary{fit,positioning}
\usepackage{pgfplots}
\pgfplotsset{compat=1.18}
\usepackage{braket}
\usepackage[utf8]{inputenc}
\usepackage{graphicx}
\usepackage{enumitem}
\usepackage{subcaption}
\usepackage{cancel}
\usepackage{multicol}
\usepackage{wrapfig}
\usepackage[T1]{fontenc}
\usepackage{url}
\usepackage{graphicx}
\usepackage{gensymb}
\usepackage{booktabs}
\usepackage{amssymb, amsmath}
\usepackage{wrapfig}
\parskip= 0.5 mm
\oddsidemargin = -1 cm
\headsep = -2 cm
\textwidth = 17.5cm
\textheight = 25.5cm
\begin{document}
\begin{center}
    {\LARGE \textbf{Inroducción a la Física Cuántica}} \\[0.4cm]

    {\large Facultad de Ciencias, UNAM} \\[0.6cm]

    \large \textbf{Tarea 4} \\[0.2cm]

    \begin{tabular}{l}
        \small Autor: Fernando Naed Ortega Montero\\
        \small Fecha: \today
    \end{tabular}
\end{center}
\[\] %borrar (Hypersnips meta.math declaration)

\section*{Preguntas conceptuales.}

\begin{enumerate}[label=\alph*)]
\item ¿Son lo mismo $ \ket{\psi} $ y $ \psi(x) $? Si no, ¿cómo se relacionan?\\

No, $ \bra{\psi} $ es un vector dentro de un espacio de Hilbert.\\
Y la funcion de onda $ \psi(x) $ esta definida como,

\[
\psi(x) := \bra{x}\ket{\psi}.
\]

Esto implica que la función de onda es una proyección del vaector ó estado del sistema ($ \ket{\psi} $) sobtre el ket de posición ($ \ket{x} $)

\item Si $ \psi(x_{0}) = 0 $ para una partícula, ¿qué puede afirmar físicamente sobre la partícula en $ x_{0} $?\\

Debido a la normalización de la función de onda, por ser considerada una distribución de probabilidad,


\[
\int |\psi(x)|^{2} \, dx = 1
\]

Esto implica que la suma de las probabilidades de todos las posibles posiciones es 1,\\
Podemos considerarslo como,

\[
\int_{x_{0}}^{x_{0}+\delta x} |\psi(x)|^{2} \, dx = 0
\]

Esto si nos da estrctamente la probabilidad de que el estado descrito por $ \psi(x) $ se encuentre entre $ x_{0} $ y $ x_{0} + \delta x$ es 0 cuando $ \delta x  \leftarrow 0$.

Asi mismo por la regla de Born la probabilida de que el sistema descrito por $ \ket{\psi} $ se enecuentre en la posición $ x_0 $ es,


\[
|\bra{x}\ket{\psi}|^{2} = |\psi(x_{0}|^{2} = 0)
\]

por lo que si $ \psi(x) = 0 $ implica que la probabilidad de que el sistema se encuentre en la psocición de $ x_{0} $ es 0.


\item Suponga que encuentra una función de onda $ \psi(x) $ que es discontinua pero cuadrado-integrable. ¿Puede descartarla inmediatamente como no física? ¿Qué otra pregunta debería hacer antes?\\

Primero habria que preguntarnos si esta función de onda es una distribución de probabilidad, osea si,

\[
\int |\psi(x)|^{2}  \, dx = 1,
\]

ó si bajo alguna constante de normalización cumple con ser una distribución de probabilidad.\\

Despues habria que considerar si la función de onda $ \psi(x) $ cumple con la ecuación de onda de Schrodinger,


\[
i\hslash \frac{\partial }{\partial t} \psi(x, t) = \left[ - \frac{\hslash^{2}}{2m} \frac{\partial^{2} }{\partial x^{2}} + \phi(x, t) \right] \psi(x, t).
\]

\end{enumerate}

\section*{Problemas}

\begin{enumerate}
\item \textbf{Su vida cambiará en 90 grados.}. Un vector $ \ket{v} $ es protectado en la dirección definida por el vector $ \ket{e_{k}} $ por la acción del proyector $ \hat{P}_{k} := \ket{e_{k}} \bra{e_{k}} $.

\begin{enumerate}
\item ¿Es la diferencia $ \ket{v} - \hat{P}_{k} \ket{v} $ ortogonal a $ \ket{e_{k}} $ ? Haga una figura de la situación y los cáluclos correspondientes para demostrar que sí o no.

Llamece al vector $ \ket{u} = \ket{v} - \hat{P}_{k} \ket{v} $, para probar que son ortogonales, el producto de vectores tiene que dar 1,


\[
\braket{u|e_{k}} = (\ket{v} - \hat{P}_{k} \ket{v})(\ket{e_{k}})
\]
\[
= (\ket{v} -  \ket{e_{k}} \braket{e_{k}|v})(\ket{e_{k}})
\]
\[
= \braket{e_{k}|v} -  \cancelto{1}{\braket{e_{k}|e_{k}}} \braket{e_{k}|v}
\]
\[
= \braket{e_{k}|v} - \braket{e_{k}|v} = 0
\]


\[
\therefore \ket{u} \mathrm{y} \ket{e_{k}} \, \mathrm{son} \, \mathrm{ortogonales}
\]

\item Sean los vectores $ \ket{v_{1}} = (1, i, 1)^t $ y $ \ket{v_{2}} = (3, 1, i)^t$. Utilice el resultado del inciso anterior para construir un par de vectores $ \ket{u_{1}} $, $ \ket{u_{2}} $ genera el mismo subespacio que $ \ket{v_{1}} $ , $ \ket{v_{2}} $. Acaba de descubrir el principio de ortonormalización de Gram-Schmidt.

Bajo la operación daga obtenemos que,


\[
\ket{v_{1}}^{\dagger} = \bra{v_{1}} = (1, -i, 1)
\]

\[
\ket{v_{2}}^{\dagger} = \bra{v_{2}} = (3, 1, -i)
\]

Considérese el producto interior de $ \ket{v_{1}} $ y $ \ket{v_{2}} $ como,

\[
\braket{v_{1}|v_{2}} = (1, -i, 1) \begin{pmatrix}
3 \\
1 \\
i \\  
\end{pmatrix} = 3
\]

Ahora difinimos un proyector tal que,

\[
\hat{P_{1}} = \ket{v_{1}} \bra{v_{1}}
\]

Y construimos el primer vector,

\[
\ket{u_{1}} = \ket{v_{2}} -\hat{P_{1}} \ket{v_{2}}
\]
\[
= \ket{v_{2}} - \ket{v_{1}}\braket{v_{1}|v_{2}}
\]
\[
= \ket{v_{2}} - 3\ket{v_{1}}
\]
\[
= \begin{pmatrix}
3 \\
1 \\
i \\  
\end{pmatrix} - 3\begin{pmatrix}
1 \\
i \\
1 \\  
\end{pmatrix}
\]
\[
\ket{u_{1}} = \begin{pmatrix}
0 \\
1-3i\\ 
-3-i \\  
\end{pmatrix} 
\]

Ahora definamos otro proyector tal que,

\[
\hat{P_{2}} = \ket{{u_{1}}} \bra{u_{1}}
\]

Y construimos el segundo vector,

\[
\ket{u_{2}} = \ket{v_{1}} - \hat{P_{2}} \ket{v_{1}}
\]

\[
= \ket{v_{1}} - \ket{{u_{1}}} \braket{u_{1}|v_{1}}
\]

Por construcción $ \ket{u_{1}} $ y $ \ket{v_{1}} $ son ortogonales, por lo que,

\[
\ket{u_{2}} = \ket{v_{1}} - \ket{{u_{1}}} \cancelto{1}{\braket{u_{1}|v_{1}}}
\]

\[
= \begin{pmatrix}
1 \\
i \\
1 \\  
\end{pmatrix} - \begin{pmatrix}
0 \\
1-3i \\
-3-i \\  
\end{pmatrix}
\]

\[
\ket{u_{2}} = \begin{pmatrix}
1 \\
-1+4i \\
4+i \\  
\end{pmatrix}
\]

Ahora normalizando los vectores $ \ket{u_{1}} $ $ \ket{u_{2}} $ se tiene,

\[
\braket{u_{1}|u_{1}} = (0, 1+3i, -3+i)\begin{pmatrix}
0 \\
1-3i\\ 
-3-i \\  
\end{pmatrix}
\]
\[
= \sqrt{0+(1-3i)(1+3i)+(-3+i)(-3-i)} = \sqrt{1+9+9+1} = \sqrt{20} 
\]

\[
\braket{u_{2}|u_{2}} = (1, -1-4i, 4-i)\begin{pmatrix}
1 \\
-1+4i \\
4+i \\  
\end{pmatrix}
\]
\[
= \sqrt{1+(-1-4i)(-1+4i)+(4-i)(4+i)} = \sqrt{1+1+16+16+1} = \sqrt{35} 
\]

\[
\ket{\hat{u_{1}}} = \frac{1}{\sqrt{20}} \ket{u_{1}}
\]
\[
\ket{\hat{u_{2}}} = \frac{1}{\sqrt{35}} \ket{u_{2}}
\]


\end{enumerate}

\end{enumerate}


\end{document}
























